\chaptertitle{Contracrostipunctus}

                               Contracrostipunctus
             Achilles has come to visit his friend and jogging companion, the
                                  Tortoise, at his home

Achilles: Heavens, you certainly have an admirable boomerang collection
Tortoise: Oh, pshaw. No better than that of any other Tortoise. And now would you like
       to step into the parlor?
Achilles: Fine. (Walks to the corner of the room.) I see you also have a large collection of
       records. What sort of music do you enjoy?
Tortoise: Sebastian Bach isn't so bad, in my opinion. But these days, I must say, I am
       developing more and more of an interest in a rather specialized sort of music.
Achilles: Tell me, what kind of music is that?
Tortoise: A type of music which you are most unlikely to have heard of. call it "music to
       break phonographs by".
Achilles: Did you say "to break phonographs by"? That is a curious concept. I can just
       see you, sledgehammer in hand, whacking on phonograph after another to pieces,
       to the strains of Beethoven's heroic masterpiece Wellington's Victory.
Tortoise: That's not quite what this music is about. However, you might find its true
       nature just as intriguing. Perhaps I should give you a brief description of it?
Achilles: Exactly what I was thinking.
Tortoise: Relatively few people are acquainted with it. It all began whet my friend the
       Crab-have you met him, by the way?-paid m• a visit.
Achilles: ' twould be a pleasure to make his acquaintance, I'm sure Though I've heard so
       much about him, I've never met him
Tortoise: Sooner or later I'll get the two of you together. You'd hit it of splendidly.
       Perhaps we could meet at random in the park on day ...
Achilles: Capital suggestion! I'll be looking forward to it. But you were going to tell me
       about your weird "music to smash phone graphs by", weren't you?
Tortoise: Oh, yes. Well, you see, the Crab came over to visit one day. You must
       understand that he's always had a weakness for fang gadgets, and at that time he
       was quite an aficionado for, of al things, record players. He had just bought his
       first record player, and being somewhat gullible, believed every word the
       salesman had told him about it-in particular, that it was capable of reproducing
       any and all sounds. In short, he was convinced that it was a Perfect phonograph.




Contracrostipunctus                                                                      83

Achilles: Naturally, I suppose you disagreed.
Tortoise: True, but he would hear nothing of my arguments. He staunchly maintained that
       any sound whatever was reproducible on his machine. Since I couldn't convince
       him of the contrary, I left it at that. But not long after that, I returned the visit,
       taking with me a record of a song which I had myself composed. The song was
       called "I Cannot Be Played on Record Player 1".
Achilles: Rather unusual. Was it a present for the Crab?
Tortoise: Absolutely. I suggested that we listen to it on his new phonograph, and he was
       very glad to oblige me. So he put it on. But unfortunately, after only a few notes,
       the record player began vibrating rather severely, and then with a loud "pop",
       broke into a large number of fairly small pieces, scattered all about the room. The
       record was utterly destroyed also, needless to say.
Achilles: Calamitous blow for the poor fellow, I'd say. What was the matter with his
       record player?
Tortoise: Really, there was nothing the matter, nothing at all. It simply couldn't reproduce
       the sounds on the record which I had brought him, because they were sounds that
       would make it vibrate and break.
Achilles: Odd, isn't it? I mean, I thought it was a Perfect phonograph. That's what the
       salesman had told him, after all.
Tortoise: Surely, Achilles, you don't believe everything that salesmen tell you! Are you
       as naive as the Crab was?
Achilles: The Crab was naiver by far! I know that salesmen are notorious prevaricators. I
       wasn't born yesterday!
Tortoise: In that case, maybe you can imagine that this particular salesman had somewhat
       exaggerated the quality of the Crab's piece of equipment ... perhaps it was indeed
       less than Perfect, and could not reproduce every possible sound.
Achilles: Perhaps that is an explanation. But there's no explanation for the amazing
       coincidence that your record had those very sounds on it ...
Tortoise: Unless they got put there deliberately. You see, before returning the Crab's
       visit, I went to the store where the Crab had bought his machine, and inquired as
       to the make. Having ascertained that, I sent off to the manufacturers for a
       description of its design. After receiving that by return mail, I analyzed the entire
       construction of the phonograph and discovered a certain set of sounds which, if
       they were produced anywhere in the vicinity, would set the device to shaking and
       eventually to falling apart.
Achilles: Nasty fellow! You needn't spell out for me the last details: that you recorded
       those sounds yourself, and offered the dastardly item as a gift ...




Contracrostipunctus                                                                       84

Tortoise: Clever devil! You jumped ahead of the story! But that wasn't t end of the
       adventure, by any means, for the Crab did r believe that his record player was at
       fault. He was quite stubborn. So he went out and bought a new record player, this
       o even more expensive, and this time the salesman promised give him double his
       money back in case the Crab found a soul which it could not reproduce exactly.
       So the Crab told r excitedly about his new model, and I promised to come over
       and see it.
Achilles: Tell me if I'm wrong-I bet that before you did so, you on again wrote the
       manufacturer, and composed and recorded new song called "I Cannot Be Played
       on Record Player based on the construction of the new model.
Tortoise: Utterly brilliant deduction, Achilles. You've quite got the spirit.
Achilles: So what happened this time?
Tortoise: As you might expect, precisely the same thing. The phonograph fell into
       innumerable pieces, and the record was shattered. Achilles: Consequently, the
       Crab finally became convinced that there could be no such thing as a Perfect
       record player.
Tortoise: Rather surprisingly, that's not quite what happened. He was sure that the next
       model up would fill the bill, and having twice the money, h e--
Achilles: Oho-I have an idea! He could have easily outwitted you, I obtaining a LOW-
       fidelity phonograph-one that was not capable of reproducing the sounds which
       would destroy it. In that way, he would avoid your trick.
Tortoise: Surely, but that would defeat the-original purpose-namely, to have a
       phonograph which could reproduce any sound whatsoever, even its own self-
       breaking sound, which is of coup impossible.
Achilles: That's true. I see the dilemma now. If any record player-si
Record Player X-is sufficiently high-fidelity, then when attempts to play the song "I
       Cannot Be Played on Record Player X", it will create just those vibrations which
       will cause to break. .. So it fails to be Perfect. And yet, the only way to g, around
       that trickery, namely for Record Player X to be c lower fidelity, even more
       directly ensures that it is not Perfect It seems that every record player is
       vulnerable to one or the other of these frailties, and hence all record players are
       defective.
Tortoise: I don't see why you call them "defective". It is simply an inherent fact about
       record players that they can't do all that you might wish them to be able to do. But
       if there is a defect anywhere, is not in THEM, but in your expectations of what
       they should b able to do! And the Crab was just full of such unrealistic
       expectations.




Contracrostipunctus                                                                      85

Achilles: Compassion for the Crab overwhelms me. High fidelity or low fidelity, he loses
       either way.
Tortoise: And so, our little game went on like *\_his for a few more rounds, and
       eventually our friend tried to become very smart. He got wind of the principle
       upon which I was basing my own records, and decided to try to outfox me. He
       wrote to the phonograph makers, and described a device of his own invention,
       which they built to specification. He called it "Record Player Omega". It was
       considerably more sophisticated than an ordinary record player.
Achilles: Let me guess how: Did it have no of cotton? Or
Tortoise: Let me tell you, instead. That will save some time. In the first place, Record
       Player Omega incorporated a television camera whose purpose it was to scan any
       record before playing it. This camera was hooked up to a small built-in computer,
       which would determine exactly the nature of the sounds, by looking at the groove-
       patterns.
Achilles: Yes, so far so good. But what could Record Player Omega do with this
       information?
Tortoise: By elaborate calculations, its little computer figured out what effects the sounds
       would have upon its phonograph. If it deduced that the sounds were such that they
       would cause the machine in its present configuration to break, then it did
       something very clever. Old Omega contained a device which could disassemble
       large parts of its phonograph subunit, and rebuild them in new ways, so that it
       could, in effect, change its own structure. If the sounds were "dangerous", a new
       configuration was chosen, one to which the sounds would pose no threat, and this
       new configuration would then be built by the rebuilding subunit, under direction
       of the little computer. Only after this rebuilding operation would Record Player
       Omega attempt to play the record.
Achilles: Aha! That must have spelled the end of your tricks. I bet you were a little
       disappointed.
Tortoise: Curious that you should think so ... I don't suppose that you know Godel's
       Incompleteness Theorem backwards and forwards, do you?
Achilles: Know WHOSE Theorem backwards and forwards? I've
heard of anything that sounds like that. I'm sure it's fascinating, but I'd rather hear more
       about "music to break records by". It's an amusing little story. Actually, I guess I
       can fill in the end. Obviously, there was no point in going on, and so you
       sheepishly admitted defeat, and that was that. Isn't that exactly it?
Tortoise: What! It's almost midnight! I'm afraid it's my bedtime. I'd love to talk some
       more, but really I am growing quite sleepy.




Contracrostipunctus                                                                      86

Achilles: As am 1. Well, 1 u be on my way. (As he reaches the door, he suddenly stops,
       and turns around.) Oh, how silly of me! I almost forgo brought you a little
       present. Here. (Hands the Tortoise a small neatly wrapped package.)
Tortoise: Really, you shouldn't have! Why, thank you very much indeed think I'll open it
       now. (Eagerly tears open the package, and ins discovers a glass goblet.) Oh, what
       an exquisite goblet! Did y know that I am quite an aficionado for, of all things, gl
       goblets?
Achilles: Didn't have the foggiest. What an agreeable coincidence!
Tortoise: Say, if you can keep a secret, I'll let you in on something: I trying to find a
       Perfect goblet: one having no defects of a sort in its shape. Wouldn't it be
       something if this goblet-h call it "G"-were the one? Tell me, where did you come
       across Goblet G?
Achilles: Sorry, but that's MY little secret. But you might like to know w its maker is.
Tortoise: Pray tell, who is it?
Achilles: Ever hear of the famous glassblower Johann Sebastian Bach? Well, he wasn't
       exactly famous for glassblowing-but he dabbled at the art as a hobby, though
       hardly a soul knows it-a: this goblet is the last piece he blew.
Tortoise: Literally his last one? My gracious. If it truly was made by Bach its value is
       inestimable. But how are you sure of its maker
Achilles: Look at the inscription on the inside-do you see where tletters `B', `A', `C', `H'
       have been etched?
Tortoise: Sure enough! What an extraordinary thing. (Gently sets Goblet G down on a
       shelf.) By the way, did you know that each of the four letters in\Bach's name is
       the name of a musical note?
Achilles:' tisn't possible, is it? After all, musical notes only go from ‘A’ through `G'.
Tortoise: Just so; in most countries, that's the case. But in Germany, Bach’s own
       homeland, the convention has always been similar, except that what we call `B',
       they call `H', and what we call `B-flat', they call `B'. For instance, we talk about
       Bach's "Mass in B Minor whereas they talk about his "H-moll Messe". Is that
       clear?
Achilles:        ... hmm ... I guess so. It's a little confusing: H is B, and B B-flat. I suppose
       his name actually constitutes a melody, then
Tortoise: Strange but true. In fact, he worked that melody subtly into or of his most
       elaborate musical pieces-namely, the final Contrapunctus in his Art of the Fugue.
       It was the last fugue Bach ever wrote. When I heard it for the first time, I had no
       idea how would end. Suddenly, without warning, it broke off. And the ... dead
       silence. I realized immediately that was where Bach died. It is an indescribably
       sad moment, and the effect it had o me was-shattering. In any case, B-A-C-H is
       the last theme c that fugue. It is hidden inside the piece. Bach didn't point it out




Contracrostipunctus                                                                           87

FIGURE 19. The last page of Bach's Art of the Fugue. In the original manuscript, in the
handwriting of Bach's son Carl Philipp Emanuel, is written: "N.B. In the course of this
fugue, at the point where the name B.A.C.H. was brought in as countersubject, the
composer died." (B-A-C-H in box.) I have let this final page of Bach's last fugue serve as
an epitaph.
[Music Printed by Donald Byrd's program "SMUT", developed at Indiana University]




Contracrostipunctus                                                                    88

Explicitly, but if you know about it, you can find it without much trouble. Ah, me-there
       are so many clever ways of hiding things in music .. .
Achilles: . . or in poems. Poets used to do very similar things, you know (though it's
       rather out of style these days). For instance, Lewis Carroll often hid words and
       names in the first letters (or characters) of the successive lines in poems he wrote.
       Poems which conceal messages that way are called "acrostics".
Tortoise: Bach, too, occasionally wrote acrostics, which isn't surprising. After all,
       counterpoint and acrostics, with their levels of hidden meaning, have quite a bit in
       common. Most acrostics, however, have only one hidden level-but there is no
       reason that one couldn't make a double-decker-an acrostic on top of an acrostic.
       Or one could make a "contracrostic"-where the initial letters, taken in reverse
       order, form a message. Heavens! There's no end to the possibilities inherent in the
       form. Moreover, it's not limited to poets; anyone could write acrostics-even a
       dialogician.
Achilles: A dial-a-logician? That's a new one on me.
Tortoise: Correction: I said "dialogician", by which I meant a writer of dialogues. Hmm
       ... something just occurred to me. In the unlikely event that a dialogician should
       write a contrapuntal acrostic in homage to J. S. Bach, do you suppose it would be
       more proper for him to acrostically embed his OWN name-or that of Bach? Oh,
       well, why worry about such frivolous matters? Anybody who wanted to write
       such a piece could make up his own mind. Now getting back to Bach's melodic
       name, did you know that the melody B-A-C-H, if played upside down and
       backwards, is exactly the same as the original?
Achilles: How can anything be played upside down? Backwards, I can see-you get H-C-
       A-B-but upside down? You must be pulling my leg.
Tortoise: ' pon my word, you're quite a skeptic, aren't you? Well, I guess I'll have to give
       you a demonstration. Let me just go and fetch my fiddle- (Walks into the next
       room, and returns in a jiffy with an ancient-looking violin.) -and play it for you
       forwards and backwards and every which way. Let's see, now ... (Places his copy
       of the Art of the Fugue on his music stand and opens it to the last page.) ... here's
       the last Contrapunctus, and here's the last theme ...

      The Tortoise begins to play: B-A-C- - but as he bows the final H, suddenly,
      without warning, a shattering sound rudely interrupts his performance. Both
      he and Achilles spin around, just in time to catch a glimpse of myriad
      fragments of glass tinkling to the floor from the shelf where Goblet G had
      stood, only moments before. And then ... dead silence.




Contracrostipunctus                                                                      89


                     Consistency, Completeness,
                           and Geometry
                           Implicit and Explicit Meaning
IN CHAPTER II, we saw how meaning-at least in the relatively simple context of formal
systems-arises when there is an isomorphism between rule-governed symbols, and things
in the real world. The more complex the isomorphism, in general, the more "equipment"-
both hardware and software-is required to extract the meaning from the symbols. If an
isomorphism is very simple (or very familiar), we are tempted to say that the meaning
which it allows us to see is explicit. We see the meaning without seeing the isomorphism.
The most blatant example is human language, where people often attribute meaning to
words in themselves, without being in the slightest aware of the very complex
"isomorphism" that imbues them with meanings. This is an easy enough error to make. It
attributes all the meaning to the object (the word), rather than to the link between that
object and the real world. You might compare it to the naive belief that noise is a
necessary side effect of any collision of two objects. This is a false belief; if two objects
collide in a vacuum, there will be no noise at all. Here again, the error stems from
attributing the noise exclusively to the collision, and not recognizing the role of the
medium, which carries it from the objects to the ear.
        Above, I used the word "isomorphism" in quotes to indicate that it must be taken
with a grain of salt. The symbolic processes which underlie the understanding of human
language are so much more complex than the symbolic processes in typical formal
systems, that, if we want to continue thinking of meaning as mediated by isomorphisms,
we shall have to adopt a far more flexible conception of what isomorphisms can be than
we have up till now. In my opinion, in fact, the key element in answering the question
"What is consciousness?" will be the unraveling of the nature of the "isomorphism"
which underlies meaning.

                  Explicit Meaning of the Contracrostipunctus
All this is by way of preparation for a discussion of the Contracrostipunctus-a study in
levels of meaning. The Dialogue has both explicit and implicit meanings. Its most
explicit meaning is simply the story




Consistency, Completeness, and Geometry                                                   90

Which was related. This “explicit meaning is, strictly speaking extremely implicit, in the
sense that the brain processes required to understand the events in the story, given only
the black marks on paper, are incredibly complex. Nevertheless, we shall consider the
events in the story to be the explicit meaning of the Dialogue, and assume that every
reader of English uses more or less the same "isomorphism" in sucking that meaning
from the marks on the paper.

         Even so, I'd like to be a little more explicit about the explicit meaning of the story.
First I'll talk about the record players and the records. The main point is that there are two
levels of meaning for the grooves in the records. Level One is that of music. Now what is
"music"-a sequence of vibrations in the air, or a succession of emotional responses in a
brain? It is both. But before there can be emotional responses, there have to be vibrations.
Now the vibrations get "pulled" out of the grooves by a record player, a relatively
straightforward device; in fact you can do it with a pin, just pulling it down the grooves.
After this stage, the ear converts the vibrations into firings of auditory neurons in the
brain. Then ensue a number of stages in the brain, which gradually transform the linear
sequence of vibrations into a complex pattern of interacting emotional responses-far too
complex for us to go into here, much though I would like to. Let us therefore content
ourselves with thinking of the sounds in the air as the "Level One" meaning of the
grooves.
What is the Level Two meaning of the grooves? It is the sequence of vibrations induced
in the record player. This meaning can only arise after the Level One meaning has been
pulled out of the grooves, since the vibrations in the air cause the vibrations in the
phonograph. Therefore, the Level Two meaning depends upon a chain of two
isomorphisms:

               (1) Isomorphism between arbitrary groove patterns and air
                  vibrations;
               (2) Isomorphism between graph vibrations. arbitrary air
                  vibrations and phonograph vibrations

This chain of two isomorphisms is depicted in Figure 20. Notice that isomorphism I is the
one which gives rise to the Level One meaning. The Level Two meaning is more implicit
than the Level One meaning, because it is mediated by the chain of two isomorphisms. It
is the Level Two meaning which "backfires", causing the record player to break apart.
What is of interest is that the production of the Level One meaning forces the production
of the Level Two meaning simultaneously-there is no way to have Level One without
Level Two. So it was the implicit meaning of the record which turned back on it, and
destroyed it.
         Similar comments apply to the goblet. One difference is that the mapping from
letters of the alphabet to musical notes is one more level of isomorphism, which we could
call "transcription". That is followed by "translation"-conversion of musical notes into
musical sounds. Thereafter, the vibrations act back on the goblet just as they did on the
escalating series of phonographs.




Consistency, Completeness, and Geometry                                                      91

FIGURE 20. Visual rendition of the principle underlying Gödel’s Theorem: two back-to-
back mappings which have an unexpected boomeranging effect. The first is from groove
patterns to sounds, carried out by a phonograph. The second-familiar, but usually
ignored -- is from sounds to vibrations of the phonograph. Note that the second mapping
exists independently of the first one, for any sound in the vicinity, not just ones produced
by the phonograph itself, will cause such vibrations. The paraphrase of Gödel’s Theorem
says that for any record player, there are records which it cannot play because they will
cause its indirect self-destruction. [Drawing by the author.



                    Implicit Meanings of the Contracrostipunctus
What about implicit meanings of the Dialogue? (Yes, it has more than one of these.) The
simplest of these has already been pointed out in the paragraphs above-namely, that the
events in the two halves of the dialogue are roughly isomorphic to each other: the
phonograph becomes a violin, the Tortoise becomes Achilles, the Crab becomes the
Tortoise, the grooves become the etched autograph, etc. Once you notice this simple
isomorphism, you can go a little further. Observe that in the first half of the story, the
Tortoise is the perpetrator of all the mischief, while in the second half, he is the victim.
What do you know, but his own method has turned around and backfired on him!
Reminiscent of the backfiring of the records' muusic-or the goblet's inscription-or perhaps
of the Tortoise's boomerang collection? Yes, indeed. The story is about backfiring on two
levels, as follows ...

       Level One: Goblets and records which backfire;
       Level Two: The Tortoise's devilish method of exploiting implicit meaning to
                  cause backfires-which backfires.

        Therefore we can even make an isomorphism between the two levels of the story,
in which we equate the way in which the records and goblet boomerang back to destroy
themselves, with the way in which the Tortoise's own fiendish method boomerangs back
to get him in the end. Seen this




Consistency, Completeness, and Geometry                                                  92

way, the story itself is an example of the backfirings which it discusses. So we can think
of the Contracrostipunctus as referring to itself indirectly that its own structure is
isomorphic to the events it portrays. (Exactly goblet and records refer implicitly to
themselves via the back-to-back morphisms of playing and vibration-causing.) One may
read the Dialogue without perceiving this fact, of course-but it is there all the time.

                      Mapping Between the Contracrostipunctus
                               and Gödel’s Theorem
        Now you may feel a little dizzy-but the best is yet to come. (Actually, levels of
implicit meaning will not even be discussed here-they will 1 for you to ferret out.) The
deepest reason for writing this Dialogue illustrate Gödel’s Theorem, which, as I said in
the Introduction, heavily on two different levels of meaning of statements of number t1
Each of the two halves of the Contracrostipunctus is an "isomorphic co Gödel’s Theorem.
Because this mapping is the central idea of the Dialogue and is rather elaborate, I have
carefully charted it out below.

                     Phonograph <= =>axiomatic system for number theory

                    low-fidelity phonograph <= =>"weak" axiomatic system

                   high-fidelity phonograph <= =>"strong" axiomatic system

               "Perfect" phonograph" <= => complete system for number theory'

               Blueprint" of phonograph <= => axioms and rules of formal system

                            record <= => string of the formal system

                     playable record<= => theorem of the axiomatic system

                  unplayable record <= =>nontheorem of the axiomatic system

                          sound <= =>true statement of number theory

                  reproducible sound <= => 'interpreted theorem of the system

                unreproducible sound <= => true statement which isn't a theorem:

                       song title <= =>implicit meaning of Gödel’s string:

                       "I Cannot Be Played            "I Cannot Be Derived
                       on Record Player X"            in Formal System X"

        This is not the full extent of the isomorphism between Gödel’s theorem and the
Contracrostipunctus, but it is the core of it. You need not if you don't fully grasp Gödel’s
Theorem by now-there are still Chapters to go before we reach it! Nevertheless, having
read this Dialogue you have already tasted some of the flavor of Gödel’s Theorem
without necessarily being aware of it. I now leave you to look for any other types of
implicit meaning in the Contracrostipunctus. "Quaerendo invenietis!"


Consistency, Completeness, and Geometry                                                  93

                                     The Art of the Fugue
A few words on the Art of the Fugue ... Composed in the last year of Bach's life, it is a
collection of eighteen fugues all based on one theme. Apparently, writing the Musical
Offering was an inspiration to Bach. He decided to compose another set of fugues on a
much simpler theme, to demonstrate the full range of possibilities inherent in the form. In
the Art of the Fugue, Bach uses a very simple theme in the most complex possible ways.
The whole work is in a single key. Most of the fugues have four voices, and they
gradually increase in complexity and depth of expression. Toward the end, they soar to
such heights of intricacy that one suspects he can no longer maintain them. Yet he does . .
. until the last Contrapunctus.
         The circumstances which caused the break-off of the Art of the Fugue (which is to
say, of Bach's life) are these: his eyesight having troubled him for years, Bach wished to
have an operation. It was done; however, it came out quite poorly, and as a consequence,
he lost his sight for the better part of the last year of his life. This did not keep him from
vigorous work on his monumental project, however. His aim was to construct a complete
exposition of fugal writing, and usage of multiple themes was one important facet of it. In
what he planned as the next-to-last fugue, he inserted his own name coded into notes as
the third theme. However, upon this very act, his health became so precarious that he was
forced to abandon work on his cherished project. In his illness, he managed to dictate to
his son-in-law a final chorale prelude, of which Bach's biographer Forkel wrote, "The
expression of pious resignation and devotion in it has always affected me whenever I
have played it; so that I can hardly say which I would rather miss-this Chorale, or the end
of the last fugue."
         One day, without warning, Bach regained his vision. But a few hours later, he
suffered a stroke; and ten days later, he died, leaving it for others to speculate on the
incompleteness of the Art of the Fugue. Could it have been caused by Bach's attainment
of self-reference?

                           Problems Caused by Gödel’s Result
        The Tortoise says that no sufficiently powerful record player can be perfect, in the
sense of being able to reproduce every possible sound from a record. Godel says that no
sufficiently powerful formal system can be perfect, in the sense of reproducing every
single true statement as a theorem. But as the Tortoise pointed out with respect to
phonographs, this fact only seems like a defect if you have unrealistic expectations of
what formal systems should be able to do. Nevertheless, mathematicians began this
century with just such unrealistic expectations, thinking that axiomatic reasoning was the
cure to all ills. They found out otherwise in 1931. The fact that truth transcends
theoremhood, in any given formal system, is called "incompleteness" of that system.
        A most puzzling fact about Gödel’s method of proof is that he uses




Consistency, Completeness, and Geometry                                                    94

reasoning methods which seemingly cannot be "encapsulated"-they re being incorporated
into any formal system. Thus, at first sight, it seems that Gödel has unearthed a hitherto
unknown, but deeply significant, difference between human reasoning and mechanical
reasoning. This mysterious discrepancy in the power of living and nonliving systems is
mirrored in the discrepancy between the notion of truth, and that of theoremhood or at
least that is a "romantic" way to view the situation.

                  The Modified pq-System and Inconsistency
        In order to see the situation more realistically, it is necessary to see in, depth why
and how meaning is mediated, in formal systems, by isomorphisms. And I believe that
this leads to a more romantic way to view i situation. So we now will proceed to
investigate some further aspects of 1 relation between meaning and form. Our first step is
to make a new formal system by modifying our old friend, the pq-system, very slightly.
We a one more axiom schema (retaining the original one, as well as the sin rule of
inference):

       Axiom SCHEMA II: If x is a hyphen-string, then xp-qx is an axiom.

Clearly, then, --p-q-- is a theorem in the new system, and so --p--q---. And yet, their
interpretations are, respectively, "2 plus; equals 2", and "2 plus 2 equals 3". It can be seen
that our new system contain a lot of false statements (if you consider strings to be
statement Thus, our new system is inconsistent with the external world.
        As if this weren't bad enough, we also have internal problems with < new system,
since it contains statements which disagree with one another such as -p-q-- (an old
axiom) and -p-q- (a new axiom). So our system is inconsistent in a second sense:
internally.
        Would, therefore, the only reasonable thing to do at this point be drop the new
system entirely? Hardly. I have deliberately presented the "inconsistencies" in a wool-
pulling manner: that is, I have tried to press fuzzy-headed arguments as strongly as
possible, with the purpose of n leading. In fact, you may well have detected the fallacies
in what I hi said. The crucial fallacy came when I unquestioningly adopted the very same
interpreting words for the new system as I had for the old of Remember that there was
only one reason for adopting those words in I last Chapter, and that reason was that the
symbols acted isomorphically to concepts which they were matched with, by the
interpretation. But when y modify the rules governing the system, you are bound to
damage t isomorphism. It just cannot be helped. Thus all the problems which we
lamented over in preceding paragraphs were bogus problems; they can made to vanish in
no time, by suitably reinterpreting some of the symbols of system. Notice that I said
"some"; not necessarily all symbols will have to mapped onto new notions. Some may
very well retain their "meaning while others change.




Consistency, Completeness, and Geometry                                                    95

Suppose, for instance, that we reinterpret just the symbol q, leaving all the others
constant; in particular, interpret q by the phrase "is greater than or equal to". Now, our
"contradictory" theorems -p-q-and -p-q--come out harmlessly as: "1 plus 1 is greater than
or equal to 1", and "1 plus 1 is greater than or equal to 2". We have simultaneously gotten
rid of (1) the inconsistency with the external world, and (2) the internal inconsistency.
And our new interpretation is a meaningful interpretation; of course the original one is
meaningless. That is, it is meaningless for the new system; for the original pq-system, it is
fine. But it now seems as pointless and arbitrary to apply it to the new pq-system as it
was to apply the "horse-apple-happy" interpretation to the old pq-system.

                        The History of Euclidean Geometry
        Although I have tried to catch you off guard and surprise you a little, this lesson
about how to interpret symbols by words may not seem terribly difficult once you have
the hang of it. In fact, it is not. And yet it is one of the deepest lessons of all of nineteenth
century mathematics! It all begins with Euclid, who, around 300 B.C., compiled and
systematized all of what was known about plane and solid geometry in his day. The
resulting work, Euclid's Elements, was so solid that it was virtually a bible of geometry
for over two thousand years-one of the most enduring works of all time. Why was this
so?
        The principal reason was that Euclid was the founder of rigor in mathematics. The
Elements began with very simple concepts, definitions, and so forth, and gradually built
up a vast body of results organized in such a way that any given result depended only on
foregoing results. Thus, there was a definite plan to the work, an architecture which made
it strong and sturdy.
        Nevertheless, the architecture was of a different type from that of, say, a
skyscraper. (See Fig. 21.) In the latter, that it is standing is proof enough that its structural
elements are holding it up. But in a book on geometry, when each proposition is claimed
to follow logically from earlier propositions, there will be no visible crash if one of the
proofs is invalid. The girders and struts are not physical, but abstract. In fact, in Euclid's
Elements, the stuff out of which proofs were constructed was human language-that
elusive, tricky medium of communication with so many hidden pitfalls. What, then, of
the architectural strength of the Elements? Is it certain that it is held up by solid structural
elements, or could it have structural weaknesses?
        Every word which we use has a meaning to us, which guides us in our use of it.
The more common the word, the more associations we have with it, and the more deeply
rooted is its meaning. Therefore, when someone gives a definition for a common word in
the hopes that we will abide by that




Consistency, Completeness, and Geometry                                                       96

      FIGURE 21. Tower of Babel, by M. C. Escher (woodcut, 1928).




Consistency, Completeness, and Geometry                             97

definition, it is a foregone conclusion that we will not do so but will instead be guided,
largely unconsciously, by what our minds find in their associative stores. I mention this
because it is the sort of problem which Euclid created in his Elements, by attempting to
give definitions of ordinary, common words such as "point", "straight line", "circle", and
so forth. How can you define something of which everyone already has a clear concept?
The only way is if you can make it clear that your word is supposed to be a technical
term, and is not to be confused with the everyday word with the same spelling. You have
to stress that the connection with the everyday word is only suggestive. Well, Euclid did
not do this, because he felt that the points and lines of his Elements were indeed the
points and lines of the real world. So by not making sure that all associations were
dispelled, Euclid was inviting readers to let their powers of association run free ...
         This sounds almost anarchic, and is a little unfair to Euclid. He did set down
axioms, or postulates, which were supposed to be used in the proofs of propositions. In
fact, nothing other than those axioms and postulates was supposed to be used. But this is
where he slipped up, for an inevitable consequence of his using ordinary words was that
some of the images conjured up by those words crept into the proofs which he created.
However, if you read proofs in the Elements, do not by any means expect to find glaring
"jumps" in the reasoning. On the contrary, they are very subtle, for Euclid was a
penetrating thinker, and would not have made any simpleminded errors. Nonetheless,
gaps are there, creating slight imperfections in a classic work. But this is not to be
complained about. One should merely gain an appreciation for the difference between
absolute rigor and relative rigor. In the long run, Euclid's lack of absolute rigor was the
cause of some of the most fertile path-breaking in mathematics, over two thousand years
after he wrote his work.
         Euclid gave five postulates to be used as the "ground story" of the infinite
skyscraper of geometry, of which his Elements constituted only the first several hundred
stories. The first four postulates are rather terse and elegant:

      (1) A straight line segment can be drawn joining any two points.

      (2) Any straight line segment can be extended indefinitely in a straight line.

      (3) Given any straight line segment, a circle can be drawn having the segment as
           radius and one end point as center.

      (4) All right angles are congruent.

The fifth, however, did not share their grace:

      (5) If two lines are drawn which intersect a third in such a way that the sum of the
            inner angles on one side is less than two right angles, then the two lines
            inevitably must intersect each other on that side if extended far enough




Consistency, Completeness, and Geometry                                                 98

Though he never explicitly said so, Euclid considered this postulate to be somehow
inferior to the others, since he managed to avoid using it in t proofs of the first twenty-
eight propositions. Thus, the first twenty-eight propositions belong to what might be
called "four-postulate geometry" that part of geometry which can be derived on the basis
of the first to postulates of the Elements, without the help of the fifth postulate. (It is al
often called absolute geometry.) Certainly Euclid would have found it 1 preferable to
prove this ugly duckling, rather than to have to assume it. B he found no proof, and
therefore adopted it.
         But the disciples of Euclid were no happier about having to assume this fifth
postulate. Over the centuries, untold numbers of people ga untold years of their lives in
attempting to prove that the fifth postulate s itself part of four-postulate geometry. By
1763, at least twenty-eight deficient proofs had been published-all erroneous! (They were
all criticized the dissertation of one G. S. Klugel.) All of these erroneous proofs involve a
confusion between everyday intuition and strictly formal properties. It safe to say that
today, hardly any of these "proofs" holds any mathematic or historical interest-but there
are certain exceptions.

                           The Many Faces of Noneuclid
        Girolamo Saccheri (1667-1733) lived around Bach's time. He had t ambition to
free Euclid of every flaw. Based on some earlier work he h; done in logic, he decided to
try a novel approach to the proof of the famous fifth: suppose you assume its opposite;
then work with that as your fif postulate ... Surely after a while you will create a
contradiction. Since i mathematical system can support a contradiction, you will have
shown t unsoundness of your own fifth postulate, and therefore the soundness Euclid's
fifth postulate. We need not go into details here. Suffice it to s that with great skill,
Saccheri worked out proposition after proposition "Saccherian geometry" and eventually
became tired of it. At one point, decided he had reached a proposition which was
"repugnant to the nature of the straight line". That was what he had been hoping for-to his
mind was the long-sought contradiction. At that point, he published his work under the
title Euclid Freed of Every Flaw, and then expired.
        But in so doing, he robbed himself of much posthumous glory, sir he had
unwittingly discovered what came later to be known as "hyperbolic geometry". Fifty
years after Saccheri, J. H. Lambert repeated the "near miss", this time coming even
closer, if possible. Finally, forty years after Lambert, and ninety years after Saccheri,
non-Euclidean geometry was recognized for what it was-an authentic new brand of
geometry, a bifurcation the hitherto single stream of mathematics. In 1823, non-
Euclidean geometry was discovered simultaneously, in one of those inexplicable
coincidences, by a Hungarian mathematician, Janos (or Johann) Bolyai, age twenty-one,
and a Russian mathematician, Nikolay Lobachevskiy, ag thirty. And, ironically, in that
same year, the great French mathematician




Consistency, Completeness, and Geometry                                                    99

Adrien-Marie Legendre came up with what he was sure was a proof of Euclid's fifth
postulate, very much along the lines of Saccheri.
        Incidentally, Bolyai's father, Farkas (or Wolfgang) Bolyai, a close friend of the
great Gauss, invested much effort in trying to prove Euclid's fifth postulate. In a letter to
his son Janos, he tried to dissuade him from thinking about such matters:

   You must not attempt this approach to parallels. I know this way to its very end. I have
   traversed this bottomless night, which extinguished all light and joy of my life. I entreat
   you, leave the science of parallels alone.... I thought I would sacrifice myself for the sake
   of the truth. I was ready to become a martyr who would remove the flaw from geometry
   and return it purified to mankind. I accomplished monstrous, enormous labors; my
   creations are far better than those of others and yet I have not achieved complete
   satisfaction. For here it is true that si paullum a summo discessit, vergit ad imum. I turned
   back when I saw that no man can reach the bottom of this night. I turned back unconsoled,
   pitying myself and all mankind.... I have traveled past all reefs of this infernal Dead Sea
   and have always come back with broken mast and torn sail. The ruin of my disposition and
   my fall date back to this time. I thoughtlessly risked my life and happiness sut Caesar aut
   nihil.'

       But later, when convinced his son really "had something", he urged him to
publish it, anticipating correctly the simultaneity which is so frequent in scientific
discovery:

   When the time is ripe for certain things, these things appear in different places in
   the manner of violets coming to light in early spring.

        How true this was in the case of non-Euclidean geometry! In Germany, Gauss
himself and a few others had more or less independently hit upon non-Euclidean ideas.
These included a lawyer, F. K. Schweikart, who in 1818 sent a page describing a new
"astral" geometry to Gauss; Schweikart's nephew, F. A. Taurinus, who did non-Euclidean
trigonometry; and F. L. Wachter, a student of Gauss, who died in 1817, aged twenty-five,
having found several deep results in non-Euclidean geometry.
        The clue to non-Euclidean geometry was "thinking straight" about the
propositions which emerge in geometries like Saccheri's and Lambert's. The Saccherian
propositions are only "repugnant to the nature of the straight line" if you cannot free
yourself of preconceived notions of what "straight line" must mean. If, however, you can
divest yourself of those preconceived images, and merely let a "straight line" be
something which satisfies the new propositions, then you have achieved a radically new
viewpoint.

                                     Undefined Terms
This should begin to sound familiar. In particular, it harks back to the pq-system, and its
variant, in which the symbols acquired passive meanings by virtue of their roles in
theorems. The symbol q is especially interesting,




Consistency, Completeness, and Geometry                                                        100

since its "meaning" changed when a new axiom schema was added. In the very same
way, one can let the meanings of "point", "line", and so on I determined by the set of
theorems (or propositions) in which they occur. This was th great realization of the
discoverers of non-Euclidean geometry. The found different sorts of non-Euclidean
geometries by denying Euclid's fifth postulate in different ways and following out the
consequences. Strict] speaking, they (and Saccheri) did not deny the fifth postulate
directly, but rather, they denied an equivalent postulate, called the parallel postulate,
which runs as follows:

   Given any straight line, and a point not on it, there exists one, and only one, straight
   line which passes through that point and never intersects the first line, no matter
   how far they are extended.

         The second straight line is then said to be parallel to the first. If you assert that no
such line exists, then you reach elliptical geometry; if you assert that, at east two such
lines exist, you reach hyperbolic geometry. Incidentally, tf reason that such variations are
still called "geometries" is that the cot element-absolute, or four-postulate, geometry-is
embedded in them. is the presence of this minimal core which makes it sensible to think
of the] as describing properties of some sort of geometrical space, even if the spa( is not
as intuitive as ordinary space.
         Actually, elliptical geometry is easily visualized. All "points", "lines and so forth
are to be parts of the surface of an ordinary sphere. Let t write "POINT" when the
technical term is meant, and "point" when t1 everyday sense is desired. Then, we can say
that a POINT consists of a pa of diametrically opposed points of the sphere's surface. A
LINE is a great circle on the sphere (a circle which, like the equator, has its center at tI
center of the sphere). Under these interpretations, the propositions ( elliptical geometry,
though they contain words like "POINT" and "LINE speak of the goings-on on a sphere,
not a plane. Notice that two LINT always intersect in exactly two antipodal points of the
sphere's surface that is, in exactly one single POINT! And just as two LINES determine
POINT, so two POINTS determine a LINE.
         By treating words such as "POINT" and "LINE" as if they had only tt meaning
instilled in them by the propositions in which they occur, we take step towards complete
formalization of geometry. This semiformal version still uses a lot of words in English
with their usual meanings (words such "the", ` if ", "and", "join", "have"), although the
everyday meaning has bee drained out of special words like "POINT" and "LINE", which
are consequently called undefined terms. Undefined terms, like the p and q of th pq-
system, do get defined in a sense: implicitly-by the totality of all propos dons in which
they occur, rather than explicitly, in a definition.
         One could maintain that a full definition of the undefined tern resides in the
postulates alone, since the propositions which follow from them are implicit in the
postulates already. This view would say that the postulates are implicit definitions of all
the undefined terms, all of the undefined terms being defined in terms of the others.




Consistency, Completeness, and Geometry                                                     101

                   The Possibility of Multiple Interpretations
A full formalization of geometry would take the drastic step of making every term
undefined-that is, turning every term into a "meaningless" symbol of a formal system. I
put quotes around "meaningless" because, as you know, the symbols automatically pick
up passive meanings in accordance with the theorems they occur in. It is another
question, though, whether people discover those meanings, for to do so requires finding a
set of concepts which can be linked by an isomorphism to the symbols in the formal
system. If one begins with the aim of formalizing geometry, presumably one has an
intended interpretation for each symbol, so that the passive meanings are built into the
system. That is what I did for p and q when I first created the pq-system.
        But there may be other passive meanings which are potentially perceptible, which
no one has yet noticed. For instance, there were the surprise interpretations of p as
"equals" and q as "taken from", in the original pq-system. Although this is rather a trivial
example, it contains the essence of the idea that symbols may have many meaningful
interpretations-it is up to the observer to look for them.
        We can summarize our observations so far in terms of the word "consistency".
We began our discussion by manufacturing what appeared to be an inconsistent formal
system-one which was internally inconsistent, as well as inconsistent with the external
world. But a moment later we took it all back, when we realized our error: that we had
chosen unfortunate interpretations for the symbols. By changing the interpretations, we
regained consistency! It now becomes clear that consistency is not a property of a formal
system per se, but depends on the interpretation which is proposed for it. By the same
token, inconsistency is not an intrinsic property of any formal system.

                              Varieties of Consistency
       We have been speaking of "consistency" and "inconsistency" all along, without
defining them. We have just relied on good old everyday notions. But now let us say
exactly what is meant by consistency of a formal system (together with an interpretation):
that every theorem, when interpreted, becomes a true statement. And we will say that
inconsistency occurs when there is at least one false statement among the interpreted
theorems.
       This definition appears to be talking about inconsistency with the external world-
what about internal inconsistencies? Presumably, a system would be internally
inconsistent if it contained two or more theorems whose interpretations were
incompatible with one another, and internally consistent if all interpreted theorems were
compatible with one another. Consider, for example, a formal system which has only the
following three theorems: TbZ, ZbE, and EbT. If T is interpreted as "the Tortoise", Z as
"Zeno", E as "Egbert", and x by as "x beats y in chess always", then we have the
following interpreted theorems:




Consistency, Completeness, and Geometry                                                 102

       The Tortoise always beats Zeno at chess
       Zeno always beats Egbert at chess.
       Egbert always beats the Tortoise at chess.

The statements are not incompatible, although they describe a rather bizarre circle of
chess players. Hence, under this interpretation, the form; system in which those three
strings are theorems is internally consistent although, in point of fact, none of the three
statements is true! Intern< consistency does not require all theorems to come out true, but
merely that they come out compatible with one another.
        Now suppose instead that x by is to be interpreted as "x was invented by y". Then
we would have:

       The Tortoise was invented by Zeno.
       Zeno was invented by Egbert.
       Egbert was invented by the Tortoise.

In this case, it doesn't matter whether the individual statements are true c false-and
perhaps there is no way to know which ones are true, and which are not. What is
nevertheless certain is that not all three can be true at one Thus, the interpretation makes
the system internally inconsistent. The internal inconsistency depends not on the
interpretations of the three capital letters, but only on that of b, and on the fact that the
three capita are cyclically permuted around the occurrences of b. Thus, one can have
internal inconsistency without having interpreted all of the symbols of the formal system.
(In this case it sufficed to interpret a single symbol.) By tl time sufficiently many symbols
have been given interpretations, it may t clear that there is no way that the rest of them
can be interpreted so that a theorems will come out true. But it is not just a question of
truth-it is question of possibility. All three theorems would come out false if the capitals
were interpreted as the names of real people-but that is not why we would call the system
internally inconsistent; our grounds for doing s would be the circularity, combined with
the interpretation of the letter I (By the way, you'll find more on this "authorship triangle"
in Chapter XX.;

                      Hypothetical Worlds and Consistency
        We have given two ways of looking at consistency: the first says that system-
plus-interpretation is consistent with the external world if every theorem comes out true
when interpreted; the second says that a system-plus: interpretation is internally
consistent if all theorems come out mutually compatible when interpreted. Now there is a
close relationship between these two types of consistency. In order to determine whether
several statements at mutually compatible, you try to imagine a world in which all of
them could be simultaneously true. Therefore, internal consistency depends upon
consistency with the external world-only now, "the external world" allowed to be any
imaginable world, instead of the one we live in. But this is




Consistency, Completeness, and Geometry                                                   103

an extremely vague, unsatisfactory conclusion. What constitutes an “imaginable" world?
After all, it is possible to imagine a world in which three characters invent each other
cyclically. Or is it? Is it possible to imagine a world in which there are square circles? Is a
world imaginable in which Newton's laws, and not relativity, hold? Is it possible to
imagine a world in which something can be simultaneously green and not green? Or a
world in which animals exist which are not made of cells? In which Bach improvised an
eight-part fugue on a theme of King Frederick the Great? In which mosquitoes are more
intelligent than people? In which tortoises can play football-or talk? A tortoise talking
football would be an anomaly, of course.
        Some of these worlds seem more imaginable than others, since some seem to
embody logical contradictions-for example, green and not green-while some of them
seem, for want of a better word, "plausible" -- such as Bach improvising an eight-part
fugue, or animals which are not made of cells. Or even, come to think of it, a world in
which the laws of physics are different ... Roughly, then, it should be possible to establish
different brands of consistency. For instance, the most lenient would be "logical
consistency", putting no restraints on things at all, except those of logic. More
specifically, a system-plus-interpretation would be logically consistent just as long as no
two of its theorems, when interpreted as statements, directly contradict each other; and
mathematically consistent just as long as interpreted theorems do not violate
mathematics; and physically consistent just as long as all its interpreted theorems are
compatible with physical law; then comes biological consistency, and so on. In a
biologically consistent system, there could be a theorem whose interpretation is the
statement "Shakespeare wrote an opera", but no theorem whose interpretation is the
statement "Cell-less animals exist". Generally speaking, these fancier kinds of
inconsistency are not studied, for the reason that they are very hard to disentangle from
one another. What kind of inconsistency, for example, should one say is involved in the
problem of the three characters who invent each other cyclically? Logical? Physical?
Biological? Literary?
        Usually, the borderline between uninteresting and interesting is drawn between
physical consistency and mathematical consistency. (Of course, it is the mathematicians
and logicians who do the drawing-hardly an impartial crew . . .) This means that the kinds
of inconsistency which "count", for formal systems, are just the logical and mathematical
kinds. According to this convention, then, we haven't yet found an interpretation which
makes the trio of theorems TbZ, ZbE, EbT inconsistent. We can do so by interpreting b
as "is bigger than". What about T and Z and E? They can be interpreted as natural
numbers-for example, Z as 0, T as 2, and E as 11. Notice that two theorems come out
true this way, one false. If, instead, we had interpreted Z as 3, there would have been two
falsehoods and only one truth. But either way, we'd have had inconsistency. In fact, the
values assigned to T, Z, and E are irrelevant, as long as it is understood that they are
restricted to natural numbers. Once again we see a case where only some of the
interpretation is needed, in order to recognize internal inconsistency.




Consistency, Completeness, and Geometry                                                   104

                Embedding of One Formal System In Another
The preceding example, in which some symbols could have interpretations while others
didn't, is reminiscent of doing geometry in natural languag4 using some words as
undefined terms. In such a case, words are divide into two classes: those whose meaning
is fixed and immutable, and, those whose meaning is to be adjusted until the system is
consistent (these are th undefined terms). Doing geometry in this way requires that
meanings have already been established for words in the first class, somewhere outside c
geometry. Those words form a rigid skeleton, giving an underlying structure to the
system; filling in that skeleton comes other material, which ca vary (Euclidean or non-
Euclidean geometry).
        Formal systems are often built up in just this type of sequential, c hierarchical,
manner. For example, Formal System I may be devised, wit rules and axioms that give
certain intended passive meanings to its symbol Then Formal System I is incorporated
fully into a larger system with more symbols-Formal System II. Since Formal System I's
axioms and rules at part of Formal System II, the passive meanings of Formal System I
symbols remain valid; they form an immutable skeleton which then plays large role in the
determination of the passive meanings of the new symbols of Formal System II. The
second system may in turn play the role of skeleton with respect to a third system, and so
on. It is also possible-an geometry is a good example of this-to have a system (e.g.,
absolute geometry) which partly pins down the passive meanings of its undefined terms,
and which can be supplemented by extra rules or axioms, which then further restrict the
passive meanings of the undefined terms. This the case with Euclidean versus non-
Euclidean geometry.

                    Layers of Stability in Visual Perception
        In a similar, hierarchical way, we acquire new knowledge, new vocabulary or
perceive unfamiliar objects. It is particularly interesting in the case understanding
drawings by Escher, such as Relativity (Fig. 22), in which there occur blatantly
impossible images. You might think that we won seek to reinterpret the picture over and
over again until we came to interpretation of its parts which was free of contradictions-
but we dot do that at all. We sit there amused and puzzled by staircases which go eve
which way, and by people going in inconsistent directions on a sing staircase. Those
staircases are "islands of certainty" upon which we base of interpretation of the overall
picture. Having once identified them, we try extend our understanding, by seeking to
establish the relationship which they bear to one another. At that stage, we encounter
trouble. But if i attempted to backtrack-that is, to question the "islands of certainty"-s
would also encounter trouble, of another sort. There's no way of backtracking and
"undeciding" that they are staircases. They are not fishes, or whip or hands-they are just
staircases. (There is, actually, one other on t-i leave all the lines of the picture totally
uninterpreted, like the "meaningless




Consistency, Completeness, and Geometry                                                 105

                  FIGURE 22. Relativity, by M. C. Escher (lithograph, 1953).

symbols" of a formal system. This ultimate escape route is an example of a "U-mode"
response-a Zen attitude towards symbolism.)
        So we are forced, by the hierarchical nature of our perceptive processes, to see
either a crazy world or just a bunch of pointless lines. A similar analysis could be made
of dozens of Escher pictures, which rely heavily upon the recognition of certain basic
forms, which are then put together in nonstandard ways; and by the time the observer
sees the paradox on a high level, it is too late-he can't go back and change his mind about
how to interpret the lower-level objects. The difference between an Escher drawing and
non-Euclidean geometry is that in the latter, comprehensible interpretations can be found
for the undefined terms, resulting in a com




Consistency, Completeness, and Geometry                                                106

prehensible total system, whereas for the former, the end result is not reconcilable with
one's conception of the world, no matter how long or stares at the pictures. Of course, one
can still manufacture hypothetic worlds, in which Escherian events can happen ... but in
such worlds, t1 laws of biology, physics, mathematics, or even logic will be violated on
or level, while simultaneously being obeyed on another, which makes the: extremely
weird worlds. (An example of this is in Waterfall (Fig. 5), whet normal gravitation
applies to the moving water, but where the nature space violates the laws of physics.)

           Is Mathematics the Same in Every Conceivable World?
We have stressed the fact, above, that internal consistency of a form; system (together
with an interpretation) requires that there be some imaginable world-that is, a world
whose only restriction is that in it, mathematics and logic should be the same as in our
world-in which all the interpreted theorems come out true. External consistency, however
consistency with the external world-requires that all theorems come of true in the real
world. Now in the special case where one wishes to create consistent formal system
whose theorems are to be interpreted as statements of mathematics, it would seem that
the difference between the two types of consistency should fade away, since, according to
what we sat above, all imaginable worlds have the same mathematics as the real world.
Thus, i every conceivable world, 1 plus 1 would have to be 2; likewise, there would have
to be infinitely many prime numbers; furthermore, in every conceivable world, all right
angles would have to be congruent; and of cours4 through any point not on a given line
there would have to be exactly on parallel line ...
        But wait a minute! That's the parallel postulate-and to assert i universality would
be a mistake, in light of what's just been said. If in all conceivable worlds the parallel
postulate-is obeyed, then we are asserting that non-Euclidean geometry is inconceivable,
which puts us back in the same mental state as Saccheri and Lambert-surely an unwise
move. But what, then, if not all of mathematics, must all conceivable worlds share?
Could it I as little as logic itself? Or is even logic suspect? Could there be worlds where
contradictions are normal parts of existence-worlds where contradictious are not
contradictions?
        Well, in some sense, by merely inventing the concept, we have shoe that such
worlds are indeed conceivable; but in a deeper sense, they are al: quite inconceivable.
(This in itself is a little contradiction.) Quite serious] however, it seems that if we want to
be able to communicate at all, we ha, to adopt some common base, and it pretty well has
to include logic. (The are belief systems which reject this point of view-it is too logical.
particular, Zen embraces contradictions and non-contradictions with equ eagerness. This
may seem inconsistent, but then being inconsistent is pa of Zen, and so ... what can one
say?)




Consistency, Completeness, and Geometry                                                    107

         Is Number Theory the Same In All Conceivable Worlds?
         If we assume that logic is part of every conceivable world (and note that we have
not defined logic, but we will in Chapters to come), is that all? Is it really conceivable
that, in some worlds, there are not infinitely many primes? Would it not seem necessary
that numbers should obey the same laws in all conceivable worlds? Or ... is the concept
"natural number" better thought of as an undefined term, like "POINT" or "LINE"? In
that case, number theory would be a bifurcated theory, like geometry: there would be
standard and nonstandard number theories. But there would have to be some counterpart
to absolute geometry: a "core" theory, an invariant ingredient of all number theories
which identified them as number theories rather than, say, theories about cocoa or rubber
or bananas. It seems to be the consensus of most modern mathematicians and
philosophers that there is such a core number theory, which ought to be included, along
with logic, in what we consider to be "conceivable worlds". This core of number theory,
the counterpart to absolute geometry-is called Peano arithmetic, and we shall formalize it
in Chapter VIII. Also, it is now well established-as a matter of fact as a direct
consequence of Gödel’s Theorem-that number theory is a bifurcated theory, with
standard and nonstandard versions. Unlike the situation in geometry, however, the
number of "brands" of number theory is infinite, which makes the situation of number
theory considerably more complex.
         For practical purposes, all number theories are the same. In other words, if bridge
building depended on number theory (which in a sense it does), the fact that there are
different number theories would not matter, since in the aspects relevant to the real world,
all number theories overlap. The same cannot be said of different geometries; for
example, the sum of the angles in a triangle is 180 degrees only in Euclidean geometry; it
is greater in elliptic geometry, less in hyperbolic. There is a story that Gauss once
attempted to measure the sum of the angles in a large triangle defined by three mountain
peaks, in order to determine, once and for all, which kind of geometry really rules our
universe. It was a hundred years later that Einstein gave a theory (general relativity)
which said that the geometry of the universe is determined by its content of matter, so
that no one geometry is intrinsic to space itself. Thus to the question, "Which geometry is
true?" nature gives an ambiguous answer not only in mathematics, but also in physics. As
for the corresponding question, "Which number theory is true?", we shall have more to
say on it after going through Gödel’s Theorem in detail.

                                     Completenes
If consistency is the minimal condition under which symbols acquire passive meanings,
then its complementary notion, completeness, is the maximal confirmation of those
passive meanings. Where consistency is the property




Consistency, Completeness, and Geometry                                                 108

way round: "Every true statement is produced by the system". Now I refine the notion
slightly. We can't mean every true statement in th world-we mean only those which
belong to the domain which we at attempting to represent in the system. Therefore,
completeness mean! "Every true statement which can be expressed in the notation of the
system is a theorem."

   Consistency: when every theorem, upon interpretation, comes out true (in some
                imaginable world).
   Completeness: when all statements which are true (in some imaginable world), and
                which can be expressed as well-formed strings of the system, are
                theorems.

         An example of a formal system which is complete on its own mode level is the
original pq-system, with the original interpretation. All true additions of two positive
integers are represented by theorems of th system. We might say this another way: "All
true additions of two positive integers are provable within the system." (Warning: When
we start using th term "provable statements" instead of "theorems", it shows that we at
beginning to blur the distinction between formal systems and their interpretations. This is
all right, provided we are very conscious of th blurring that is taking place, and provided
that we remember that multiple interpretations are sometimes possible.) The pq-system
with the origin interpretation is complete; it is also consistent, since no false statement is-,
use our new phrase-provable within the system.
         Someone might argue that the system is incomplete, on the grounds that additions
of three positive integers (such as 2 + 3 + 4 =9) are not represented by theorems of the
pq-system, despite being translatable into the notation of the system (e.g., --p---p----q----
--------). However, this string is not well-formed, and hence should be considered to I just
as devoid of meaning as is p q p---q p q. Triple additions are simply not expressible in
the notation of the system-so the completeness of the system is preserved.
         Despite the completeness of the pq-system under this interpretation, certainly falls
far short of capturing the full notion of truth in numb theory. For example, there is no
way that the pq-system tells us how mat prime numbers there are. Gödel’s
Incompleteness Theorem says that any system which is "sufficiently powerful" is, by
virtue of its power, incomplete, in the sense that there are well-formed strings which
express tr statements of number theory, but which are not theorems. (There a truths
belonging to number theory which are not provable within the system.) Systems like the
pq-system, which are complete but not very powerful, are more like low-fidelity
phonographs; they are so poor to beg with that it is obvious that they cannot do what we
would wish them do-namely tell us everything about number theory.




Consistency, Completeness, and Geometry                                                    109

         How an Interpretation May Make or Break Completeness
What does it mean to say, as I did above, that "completeness is the maximal confirmation
of passive meanings"? It means that if a system is consistent but incomplete, there is a
mismatch between the symbols and their interpretations. The system does not have the
power to justify being interpreted that way. Sometimes, if the interpretations are
"trimmed" a little, the system can become complete. To illustrate this idea, let's look at
the modified pq-system (including Axiom Schema II) and the interpretation we used for
it.
        After modifying the pq-system, we modified the interpretation for q from "equals"
to "is greater than or equal to". We saw that the modified pq-system was consistent under
this interpretation; yet something about the new interpretation is not very satisfying. The
problem is simple: there are now many expressible truths which are not theorems. For
instance, "2 plus 3 is greater than or equal to 1" is expressed by the nontheorem --p---q-.
The interpretation is just too sloppy! It doesn't accurately reflect what the theorems in the
system do. Under this sloppy interpretation, the pq-system is not complete. We could
repair the situation either by (1) adding new rules to the system, making it more
powerful, or by (2) tightening up the interpretation. In this case, the sensible alternative
seems to be to tighten the interpretation. Instead of interpreting q as "is greater than or
equal to", we should say "equals or exceeds by 1". Now the modified pq-system becomes
both consistent and complete. And the completeness confirms the appropriateness of the
interpretation.

                Incompleteness of Formalized Number Theory
In number theory, we will encounter incompleteness again; but there, to remedy the
situation, we will be pulled in the other direction-towards adding new rules, to make the
system more powerful. The irony is that we think, each time we add a new rule, that we
surely have made the system complete now! The nature of the dilemma can be illustrated'
by the following allegory ...
        We have a record player, and we also have a record tentatively labeled "Canon on
B-A-C-H". However, when we play the record on the record player, the feedback-
induced vibrations (as caused by the Tortoise's records) interfere so much that we do not
even recognize the tune. We conclude that something is defective-either our record, or
our record player. In order to test our record, we would have to play it on friends' record
players, and listen to its quality. In order to test our phonograph, we would have to play
friends' records on it, and see if the music we hear agrees with the labels. If our record
player passes its test, then we will say the record was defective; contrariwise, if the
record passes its test, then we will say our record player was defective. What, however,
can we conclude when we find out that both pass their respective tests? That is the
moment to remember the chain of two isomorphisms (Fig. 20), and think carefully!




Consistency, Completeness, and Geometry                                                  110

